\section{Formeln zur $\Gamma$-Funktion\label{gamma:section:formulas}}
\kopfrechts{Formlel zur $\Gamma$-Funktion}

Für die $\Gamma$-Funktion existieren verschiedene Darstellungen und Näherungsverfahren.
Im Folgenden werden ausgewählte Formeln sowie ihre jeweiligen Anwendungsgebiete vorgestellt, wobei auf ausführliche mathematische Herleitungen verzichtet wird.
Abschliessend wird die Bedeutung der $\Gamma$-Funktion für die Berechnung des Kugelnvolumen aufgezeigt.

\subsection{Produktformel von Euler}

Auf der Suche nach der $\Gamma$-Funktion im 18. Jahrhundert hat Leonhard Euler unendliche Produkte verwendet. Die Formel
\begin{equation}
\label{gamma-eq:product-formula-euler}
    \Gamma(z) = \lim_{n \to \infty} \frac{n!\,n^z}{z(z+1)\cdots(z+n)}
\end{equation}
liefert die Funktionalgleichung $\Gamma(z+1)=z\Gamma(z)$. 
Zusammen mit $\Gamma(1)=1$ folgt daraus für jede natürliche Zahl $n$, dass $\Gamma(n+1)=n!$ gilt.
Dementsprechend setzt die $\Gamma$-Funktion die Fakultät auf nicht-ganzzahlige Argumente fort.

\subsection{Integraldarstellung}

Die Integraldarstellung
\begin{equation}
\label{gamma-eq:integral}
    \Gamma(z) = \int_0^\infty t^{z-1} e^{-t}\,dt,
    \qquad \operatorname{Re}(z)>0
\end{equation}
ist die Standarddefinition der $\Gamma$-Funktion. 
Sie eignet sich besonders gut zur Untersuchung von Ableitungen. 
Durch Differentiation unter dem Integralzeichen erhält man
\begin{equation}
    \Gamma'(z)=\int_0^\infty t^{z-1}e^{-t}\log(t)\,dt.
\end{equation}
Außerdem lassen sich die Werte der $\Gamma$-Funktion mithilfe numerischer Integration mit hoher Genauigkeit approximieren.
Die Produktformel von Euler aus Gleichung \eqref{gamma-eq:product-formula-euler} und die Integraldarstellung \eqref{gamma-eq:integral} definieren dieselbe Funktion.

\subsection{Lanczos-Approximation}

Die Lanczos-Approximation
\begin{equation}
\label{gamma-eq:lanczos}
    \Gamma(z) \approx \sqrt{2\pi}
    \left(z + g - \frac{1}{2}\right)^{z - \frac{1}{2}}
    e^{-\left(z + g - \frac{1}{2}\right)}
    \left(c_0 + \sum_{k=1}^{N} \frac{c_k}{z + k - 1}\right)
\end{equation}
wurde 1964 von Cornelius Lanczos vorgestellt und wird heute häufig in numerischen Bibliotheken eingesetzt.
Der Parameter $g$ und die Koeffizienten $c_k$ werden dabei im Voraus festgelegt und müssen nicht bei jeder Auswertung neu berechnet werden.
Da nur eine endliche Summe mit $N$ Gliedern ausgewertet wird, ermöglicht das Verfahren eine genaue und zugleich effiziente Berechnung der $\Gamma$-Funktion.